%% GENERATED by python/paper/build.py from results/ -- do not edit by hand.
%% Rebuild:  .venv\Scripts\python.exe python\paper\build.py --only OECD_LeisureCalibration
%% Source:   results/paper/usCalibrationSummary.csv
\begin{table}[!htb]
\centering
\begin{threeparttable}
\caption{The U.S., the UK and France under the two treatments of the taste for leisure}
\label{table:OA:OECD_LeisureCalibration}
\renewcommand{\arraystretch}{1.25}
\begin{tabularx}{.9\textwidth}{lYYYYYY}
\toprule
 & \multicolumn{2}{c}{\textbf{U.S.}} & \multicolumn{2}{c}{\textbf{UK}} & \multicolumn{2}{c}{\textbf{France}} \\
\cmidrule(lr){2-3}\cmidrule(lr){4-5}\cmidrule(lr){6-7}
\textbf{Parameter} & \textbf{Common} & \textbf{Vector} & \textbf{Common} & \textbf{Vector} & \textbf{Common} & \textbf{Vector} \\
\midrule 
$X$ & 4.07 & 0.88 & 7.67 & 1.45 & 6.48 & 1.23 \\ % row: X
$\eta_H/\eta_L$ & 3.08 & 3.73 & 3.06 & 3.87 & 2.29 & 2.41 \\ % row: etaratio
\bottomrule
\end{tabularx}
\begin{tablenotes}[flushleft]
\footnotesize
\item[] \textit{Note:} $\rho = 1$. Under the vector $X_i$, $X$ is the population-weighted mean of the $X_i$. The other parameters of table \ref{table:US:CalibUSUKFR} are the same under the two calibrations.
\end{tablenotes}
\end{threeparttable}
\end{table}
